\documentclass[12pt]{article}
\usepackage[T1]{fontenc}
\usepackage[utf8]{inputenc}
\usepackage{lmodern}
\usepackage{textcomp}
\usepackage{enumitem}
\usepackage{geometry}
\geometry{margin=1in}
\begin{document}
\begin{center}
Answer Key - Astronomy - Graduate Level Assessment 17 \
\small (Answers are ordered exactly as in the question paper.)
\end{center}

\section*{Multiple Choice}
\begin{enumerate}[label=\arabic*.]
\item \textbf{Answer:} B
\\ \textit{Options:} A) The resonance pulls Io in different directions and generates heat.; B) It makes the orbit of Io slightly elliptical.; C) It creates a gap with no asteriods between the orbits.; D) It prevents formation of the ring material into other moons.
\\ \textit{Note:} Correct option: B - It makes the orbit of Io slightly elliptical.
\item \textbf{Answer:} D
\\ \textit{Options:} A) because the Moon has a nearly circular orbit around the Earth; B) because the Moon does not rotate; C) because the other face points toward us only at new moon when we can't see the Moon; D) because the Moon's rotational and orbital periods are equal
\\ \textit{Note:} Correct option: D - because the Moon's rotational and orbital periods are equal
\item \textbf{Answer:} B
\\ \textit{Options:} A) Spontaneous emission from radioactive atoms.; B) The minimization of gravitational potential energy.; C) Thermally induced collisions.; D) Plate tectonics.
\\ \textit{Note:} Correct option: B - The minimization of gravitational potential energy.
\item \textbf{Answer:} D
\\ \textit{Options:} A) fast rotation only; B) a rocky mantle only; C) a molten metallic core only; D) both a molten metallic core and reasonably fast rotation
\\ \textit{Note:} Correct option: D - both a molten metallic core and reasonably fast rotation
\item \textbf{Answer:} D
\\ \textit{Options:} A) hydrogen compounds (H 2 O CH 4 NH 3) light gases (H He) metals rocks; B) hydrogen compounds (H 2 O CH 4 NH 3) light gases (H He) rocks metals; C) light gases (H He) hydrogen compounds (H 2 O CH 4 NH 3) metals rocks; D) light gases (H He) hydrogen compounds (H 2 O CH 4 NH 3) rocks metals
\\ \textit{Note:} Correct option: D - light gases (H He) hydrogen compounds (H 2 O CH 4 NH 3) rocks metals
\end{enumerate}

\section*{True / False}
\begin{enumerate}[label=\arabic*.]
\setcounter{enumi}{5}
\item \textbf{Answer:} False
\\ \textit{Options:} A) True; B) False
\\ \textit{Note:} LLM-generated false statement. Correct answer: 'Sagittarius A*'.
\item \textbf{Answer:} True
\\ \textit{Options:} A) True; B) False
\\ \textit{Note:} LLM-generated true statement based on MCQ.
\item \textbf{Answer:} False
\\ \textit{Options:} A) True; B) False
\\ \textit{Note:} LLM-generated false statement. Correct answer: 'Radioactive dating of lunar samples shows that they are older.'.
\item \textbf{Answer:} False
\\ \textit{Options:} A) True; B) False
\\ \textit{Note:} LLM-generated false statement. Correct answer: 'third quarter.'.
\item \textbf{Answer:} True
\\ \textit{Options:} A) True; B) False
\\ \textit{Note:} LLM-generated true statement based on MCQ.
\end{enumerate}

\section*{Long Form Response}
\begin{enumerate}[label=\arabic*.]
\setcounter{enumi}{10}
\item \textbf{Answer:} Jupiter. Explain why this is the correct answer and provide supporting evidence. Reference key facts or concepts that support this choice.
\\ \textit{Note:} Long-form derived from MCQ; award credit for stating the correct idea and giving supporting reasoning.
\item \textbf{Answer:} 1.5 x 1017 meters. Explain why this is the correct answer and provide supporting evidence. Reference key facts or concepts that support this choice.
\\ \textit{Note:} Long-form derived from MCQ; award credit for stating the correct idea and giving supporting reasoning.
\end{enumerate}

\vfill
\noindent\textit{End of Answer Key}
\end{document}